% unidades/06-gustos-preferencias.tex
\chapter{❤️ 好き・嫌い — Gustos y Preferencias}
\unidadheader{06}{好き・嫌い}{Gustos y Preferencias}{Expresar lo que te gusta y no te gusta}

\begin{tcolorbox}[vocabbox, title={📌 Vocabulario Nuevo}]
\begin{tabularx}{\linewidth}{llX}
\toprule
\textbf{Japonés} & \textbf{Español} & \textbf{Tipo} \\
\midrule
\vocabitem{好き}{すき}{gustar / que se gusta}{adj-な}
\vocabitem{嫌い}{きらい}{disgustar / que no se gusta}{adj-な}
\vocabitem{得意}{とくい}{ser bueno en ~}{adj-な}
\vocabitem{苦手}{にがて}{ser malo en ~ / no ser lo tuyo}{adj-な}
\vocabitem{趣味}{しゅみ}{hobby / pasatiempo}{sust.}
\vocabitem{音楽}{おんがく}{música}{sust.}
\vocabitem{映画}{えいが}{película / cine}{sust.}
\vocabitem{料理}{りょうり}{cocina / cocinar}{sust./v.}
\vocabitem{読書}{どくしょ}{lectura}{sust.}
\vocabitem{スポーツ}{—}{deporte}{sust.}
\bottomrule
\end{tabularx}
\end{tcolorbox}

\subsection*{🗣️ Frases Clave}

\frase{\ruby{趣味}{しゅみ}は\ruby{何}{なん}ですか？}{¿Cuál es tu hobby?}
\frase{\ruby{音楽}{おんがく}が\ruby{大好}{だいす}きです。}{Me encanta la música.}
\frase{スポーツは\ruby{苦手}{にがて}です。}{El deporte no es lo mío.}
\frase{〜はどうですか？}{¿Qué te parece ~? / ¿Cómo te va con ~?}

\subsection*{⚙️ Gramática}

\patron{Me gusta / No me gusta}
  {〜が 好き / 嫌い です}
  {好き・嫌い son adjetivos-な, no verbos.
  \ej{\ruby{私}{わたし}はラーメンが\ruby{好}{す}きです.}{Me gusta el ramen.}
  \ej{\ruby{私}{わたし}は\ruby{虫}{むし}が\ruby{嫌}{きら}いです。}{Los insectos me disgustan.}
  \ej{\ruby{猫}{ねこ}が\ruby{大好}{だいす}きです！}{¡Adoro los gatos!}}

\patron{Ser bueno / malo en algo}
  {〜が 得意 / 苦手 です}
  {\ej{\ruby{数学}{すうがく}が\ruby{得意}{とくい}です。}{Soy bueno en matemáticas.}
  \ej{\ruby{歌}{うた}が\ruby{苦手}{にがて}です。}{No soy bueno cantando.}}

\patron{¿Qué te parece…?}
  {〜はどうですか？}
  {\ej{この\ruby{映画}{えいが}はどうですか？}{¿Qué te parece esta película?}
  \ej{とても\ruby{面白}{おもしろ}いです！}{¡Es muy interesante!}}

\begin{tcolorbox}[ejerciciobox, title={📝 Ejercicios Rápidos}]
\begin{enumerate}
  \item Di dos cosas que te gustan mucho usando 大好き.
  \item Pregunta a alguien cuál es su hobby.
  \item Completa: \ruby{私}{わたし}は\ruby{料理}{りょうり}が\_\_\_\_です。(Soy bueno/a cocinando)
\end{enumerate}
\vspace{0.4em}
\textbf{Respuestas:}
1) (personal)\quad
2) \ruby{趣味}{しゅみ}は\ruby{何}{なん}ですか？\quad
3) \ruby{得意}{とくい}
\end{tcolorbox}

\begin{tcolorbox}[kanjibox, title={🀄 Kanjis de la Unidad}]
\begin{tabularx}{\linewidth}{cllXX}
\toprule
\textbf{Kanji} & \textbf{On} & \textbf{Kun} & \textbf{Significado} & \textbf{Vocab} \\
\midrule
\kanjirow{好}{こう}{す・この}{gustar}{\ruby{好き}{すき} / \ruby{好物}{こうぶつ}}
\kanjirow{嫌}{けん・けん}{きら・いや}{desagradar}{\ruby{嫌い}{きらい} / \ruby{嫌}{いや}}
\kanjirow{音}{おん・いん}{おと}{sonido / música}{\ruby{音楽}{おんがく} / \ruby{音}{おと}}
\kanjirow{楽}{がく・らく}{たの}{música / fácil / placer}{\ruby{音楽}{おんがく} / \ruby{楽しい}{たのしい}}     
\kanjirow{映}{えい}{うつ}{reflejar / proyectar}{\ruby{映画}{えいが} / \ruby{映}{うつ}る}
\bottomrule
\end{tabularx}
\end{tcolorbox}
